\begin{center}
  \large\textbf{ABSTRAK}
\end{center}

\addcontentsline{toc}{chapter}{ABSTRAK}

\vspace{2ex}

\begingroup
% Menghilangkan padding
\setlength{\tabcolsep}{0pt}

\noindent
\begin{tabularx}{\textwidth}{l >{\centering}m{2em} X}
  Nama Mahasiswa    & : & \name{}         \\

  Judul Tugas Akhir & : & \tatitle{}      \\

  Pembimbing        & : & 1. \advisor{}   \\
                    &   & 2. \coadvisor{} \\
\end{tabularx} 
\endgroup

Penyu laut merupakan satwa yang dilindungi karena populasinya terus mengalami penurunan akibat aktivitas manusia seperti perburuan ilegal dan kerusakan habitat. Upaya konservasi memerlukan sistem identifikasi individu yang akurat untuk mendukung pemantauan populasi dan pola migrasi. Metode konvensional seperti satellite tagging memiliki keterbatasan berupa biaya tinggi, risiko kerusakan perangkat, serta sifat yang invasif terhadap hewan. Oleh karena itu, penelitian ini mengembangkan model re-identifikasi (Re-ID) penyu berbasis citra kepala menggunakan arsitektur Vision Transformer (ViT) sebagai alternatif yang lebih efisien dan non-invasif.
Penelitian menggunakan dataset SeaTurtleIDHeads yang berisi citra kepala penyu dari berbagai individu dan kondisi pengambilan gambar. Model dibangun menggunakan backbone Vision Transformer berbasis pra-pelatihan DINOv3 dan dilengkapi dengan embedding head untuk menghasilkan representasi fitur visual. Proses pelatihan menerapkan kombinasi ArcFace Loss dan Online Hard Triplet Loss untuk meningkatkan separabilitas antar identitas. Selain itu, digunakan teknik Layer-wise Learning Rate Decay (LLRD), cosine learning rate scheduler, serta test-time augmentation guna meningkatkan performa model. Evaluasi dilakukan menggunakan metrik Rank-1 dan mean Average Precision (mAP) pada skenario open-set re-identification.

Kata Kunci: Penyu Laut, Re-identifikasi, Vision Transformer
